\refstepcounter{section}
\section*{Lecture 12: Making things invariant}
\addcontentsline{toc}{section}{Lecture 12: Making things invariant}
The states that we had defined earlier, $\ket{\vb{p}_1, \vb{p}_2, \ldots}$ are \textit{improper}, in the sense that they cannot be normalised. The obvious choice, $\braket{\vb{p}|\vb{p'}} = (2\pi)^3 \delt{3}{\vb{p}-\vb{p}'}$ is not physical due to the presence of the Dirac-delta. We can form physical states by `smearing out' these deltas, by defining:
$$\ket{\psi} = \int \meas{\vb{p}}\psi(\vb{p})\ket{\vb{p}} = \int \meas{\vb{p}}\psi(\vb{p})\crep\ket{\vac}$$ where $\psi(\vb{p})$ is an $\mathbb{L}^2$ function. However, this smearing process is not apriori Lorentz invariant. Thus, we can force the process to follow Lorentz invariance since the inception only, so that we do not run into trouble later.
% Let us now see how we can normalise the improper vectors $\{\ket{\vb{p}}\}$ to some new set of improper vectors, which are manifestly Lorentz invariant. \\[0.2cm]
% Define the \textit{projector} onto the single-particle state,
% $$\mathds{1} := \int \meas{\vb{p}} \ket{\vb{p}}\bra{\vb{p}}$$
% The wavefunction can change depending on which frame we are, however, whether we have a single particle or not, is frame independent. Separately, the measure $ \meas{\vb{p}}$ is not invariant and hence $\ket{\vb{p}}\bra{\vb{p}}$ is not too, however, together it produces something invariant. Let us now consider a mystery factor $X(\vb{p})$ such that $\int \meas{\vb{p}} \frac{1}{X(\vb{p})}$ is Lorentz invariant.\\[0.2cm]
% Note that $\int \dd^4 p$ is Lorentz invariant. 
\subsection{Invariant Inner Product}
For simplicity, consider a Lorentz boost along the $x$ direction. Then, 
$${p'}^1 = \gamma(p^1+\beta E)\qquad\qquad E' = \gamma(E+\beta p^1)$$
where $\gamma = \sfrac{1}{\sqrt{1-\qty(\sfrac{v}{c})^2}}$ and $\beta = \sfrac{v}{c}$ are defined as usual. Now, note that $p^\mu p_\mu = m^2 \implies E^2 = (p^0)^2 +(p^1)^2+(p^2)^2+m^2 \implies 0<E= \sqrt{(p^0)^2 +(p^1)^2+(p^2)^2+m^2}$ under the Minkowski metric. As the transformation is only along $x$ direction, the momentum in $x$ and $y$ do not change. Thus, we have:
$$\pdv{E}{p^j} = \frac{1}{2E}\times (2 p^j) \implies \pdv{E}{p^1} = \frac{p^1}{E}$$
Now, note the transformation of the Dirac-delta:
\begin{align*}
   \int \dd^3 \vb{y} \ \delt{3}{\vb{x}-\vb{y}} =  \int \dd^3 \vb{y}' \ \delt{3}{\vb{x}'-\vb{y}'} = \int \Big|\pdv{\vb{y}'}{\vb{y}}\Big| \dd^3 \vb{y} \ \delt{3}{\vb{x}'-\vb{y}'}
\end{align*}
Comparing, we see that the Dirac-delta transforms as $\delt{3}{\vb{x}-\vb{y}} = \Big|\pdv{\vb{y}'}{\vb{y}}\Big| \delt{3}{\vb{x}'-\vb{y}'}$ where $\Big|\pdv{\vb{y}'}{\vb{y}}\Big|$ is the determinant of the Jacobian matrix $J$ of the transformation. Hence, for the Lorentz transformation, we have:
\begin{align*}
   \delt{3}{\vb{p}-\vb{k}} &=  \delta(p^1-k^1) \delta(p^2-k^2) \delta(p^3-k^3)\\
   &= \mathfrak{J} \ \delta(p'^1-k'^1) \delta(p'^2-k'^2) \delta(p'^3-k'^3)
\end{align*}
Now, let us calculate the Jacobian matrix, whose elements will be given by $J^{ij} = \pdv{p'^i}{p^j}$ (note that since we are concerned only with 3-momentum delta function, only spatial indices will appear). Note that $\pdv{p'^2}{p^j} = \delta_{2j}$ and $\pdv{p'^3}{p^j} = \delta_{3j}$. Only for $j=1$, we have:
$$\pdv{p'^1}{p^j} = \gamma \delta_{1j} + \gamma\beta \pdv{E}{p^j}$$ and from the dispersion relation, we have $\pdv{E}{p^j} = \frac{p^j}{E}$. The matrix is then:
$$J \equiv \mqty(\gamma \qty(1+\beta \frac{p^1}{E}) & \gamma\beta \frac{p^2}{E} & \gamma\beta \frac{p^3}{E}\\ 0 & 1 & 0 \\ 0 & 0 & 1)\implies \mathfrak{J} =  \det J  = \gamma \qty(1+\beta \frac{p^1}{E}) = \frac{\gamma}{E} \qty(E+\beta p^1) = \frac{E'}{E}$$
Thus, we obtain the transformation of the delta function as:
$$   \delt{3}{\vb{p}-\vb{k}}= \frac{E'}{E}   \delt{3}{\vb{p}'-\vb{k}'}\implies E \delt{3}{\vb{p}-\vb{k}} \text{ is invariant under L.T} $$
Then, we can now define the Lorentz-invariant inner product such that:
$$\braket{\vb{p}|\vb{k}} =(2E_{\vb{p}}) (2\pi)^3 \delt{3}{\vb{p}-\vb{k}}\equiv (2\omp) (2\pi)^3 \delt{3}{\vb{p}-\vb{k}}$$
In line with this, we can therefore define single-particle state as:
$$\ket{\vb{p}} = \sqrt{2\omp} \crep \vac$$
\subsection{Invariant Measure}
